%-----------------------------------------------------------------------------------------------
\makeatletter
\immediate\write18{datelog > \jobname.info} % site script for $(date '+%Y-%m-%d %Hh%Mm%Ss')
\makeatother
%-----------------------------------------------------------------------------------------------
%-----------------------------------------------------------------------------------------------
\usetheme{Copenhagen}
\usepackage{beamercolorthemeCNThermSci2}
\usefonttheme{serif}
%-----------------------------------------------------------------------------------------------

%-----------------------------------------------------------------------------------------------
%-----------------------------------------------------------------------------------------------
\usetheme{Copenhagen}
\usepackage{beamercolorthemeUTF2}
\usefonttheme{serif}
%-----------------------------------------------------------------------------------------------
\usepackage[utf8]{inputenc}
\usepackage[greek,french,english,brazil]{babel} % last becomes the active one
\usepackage{pslatex}
\usepackage{amssymb,amsmath}
\usepackage{soul}
\usepackage[squaren,Gray,cdot]{SIunits}
\usepackage[nice]{nicefrac}
\usepackage{tikz}
\usepackage{amscd}
\usepackage{stmaryrd}
\usepackage{scalerel}
\usepackage{xspace}
%-----------------------------------------------------------------------------------------------


%-----------------------------------------------------------------------------------------------
%-----------------------------------------------------------------------------------------------
% Mathematical
%-----------------------------------------------------------------------------------------------
\newcommand{\vet}[1]{\underline{{#1}}}
\newcommand{\mat}[1]{\underline{\underline{{#1}}}}
\newcommand{\cub}[1]{\underline{\underline{\underline{{#1}}}}}
\newcommand{\eqdef}{{\ensuremath\stackrel{\text{\tiny def}}{=}}}
%-----------------------------------------------------------------------------------------------
% Linguistic
%-----------------------------------------------------------------------------------------------
\newcommand{\GRtxt}[1]{\begin{otherlanguage}{greek}{{#1}}\end{otherlanguage}}
\newcommand{\FRtxt}[1]{\begin{otherlanguage}{french}{{#1}}\end{otherlanguage}}
%-----------------------------------------------------------------------------------------------
% Presentation
%-----------------------------------------------------------------------------------------------
\newcommand{\BkgImgH}[1]{% Places an image centered on the slide background filling the height
    \usebackgroundtemplate{\parbox{\paperwidth}{%
        \vspace*{1sp}\centering\includegraphics[height=\paperheight]{{#1}}
}}}
\newcommand{\BkgImgW}[1]{% Places an image centered on the slide background filling the width
    \usebackgroundtemplate{\parbox{\paperwidth}{%
        \vspace*{1sp}\centering\includegraphics[width=\paperwidth]{{#1}}
}}}
\newcommand{\ArtEndH}[3]{% Transitions to plain image (last) slide: #1:prefix #2,#3:extensions
    \BkgImgH{root/../art/#1.#2}
    \frame<handout:0>[plain]{%
        \transdissolve\vspace*{72mm}\color{white}\scriptsize\bf\input{root/../art/#1.#3}}
    \usebackgroundtemplate{\mbox{~}}
}
\newcommand{\ArtEndW}[3]{% Transitions to plain image (last) slide: #1:prefix #2,#3:extensions
    \BkgImgW{root/../art/#1.#2}
    \frame<handout:0>[plain]{%
        \transdissolve\vspace*{72mm}\color{white}\scriptsize\bf\input{root/../art/#1.#3}}
    \usebackgroundtemplate{\mbox{~}}
}
\newcommand{\ImgColW}[3]{% Inserts a full-width image in a column
    \includegraphics[width=\columnwidth]{root/../art/#1.#2}\\[-0.5\baselineskip]
    \parbox{\columnwidth}{\tiny\hfill\scalebox{0.85}{\input{root/../art/#1.#3}}}
}
\newcommand{\txtpic}[1]{%
    \fcolorbox{lightgray}{white!90!black}{{#1}} 
}
%-----------------------------------------------------------------------------------------------


%-----------------------------------------------------------------------------------------------
\title{B.04.01 -- Ciclos de Refrigeração}
\subtitle{Ciclos de Compressão de Vapor de Simples Estágio}
\author{Prof.~C.~Naaktgeboren, PhD}
\date{{\scriptsize\tt%
    \includegraphics[height=6.0mm]{cc/by-nc-nd-88x31.pdf}\\[\smallskipamount]
    https://github.com/CNThermSci/ApplThermSci\\
    Compiled on \input{\jobname.info}
}}
%-----------------------------------------------------------------------------------------------
\begin{document}
%-----------------------------------------------------------------------------------------------
\logo{%
    \parbox{158mm}{% There's a 1mm gap on each side of the 160mm x 90mm slide logo line
        \mode<beamer>{
            \includegraphics[height=6.0mm]{root/00-res/UTFPR/UTFPR-logo-D.pdf}\hfill%
            \includegraphics[height=9.0mm]{root/00-res/logo/CNThermSci-logo-A.pdf}%
        }
        \mode<handout>{
            \includegraphics[height=6.0mm]{root/00-res/UTFPR/UTFPR-logo-W.pdf}\hfill%
            \includegraphics[height=9.0mm]{root/00-res/logo/CNThermSci-logo-W.pdf}%
        }
    }
} % The (delineated, alpha), or washed-out logos
%-----------------------------------------------------------------------------------------------
\frame{\titlepage}
%-----------------------------------------------------------------------------------------------

%-----------------------------------------------------------------------------------------------
\frame{\tableofcontents}
%-----------------------------------------------------------------------------------------------

%-----------------------------------------------------------------------------------------------
\frame{%
    % !bib LST 2>/dev/null | paste -s -d',' | sed 's|,|, |g' | j -i8 96
    \hspace{0.10\textwidth}
    \begin{minipage}{0.80\textwidth}
        Esta apresentação baseia-se nas seguintes referências / tópicos de leitura:
        \\[\medskipamount]
        \begin{itemize}
            \item \cite{2016-FentonDL-ASHRAE}, \alert{Capítulo~3, pp.~45--50} e
                \\[\bigskipamount]
            \item \cite{2013-CengelYA+BolesMA-AMGH}, \alert{Seções~11--1 e 11--3 a 7}.
        \end{itemize}
    \end{minipage}
}
%-----------------------------------------------------------------------------------------------

%-----------------------------------------------------------------------------------------------
\section{Ciclos de Refrigeração por Compressão de Vapor Simples Estágio}
%-----------------------------------------------------------------------------------------------

    %-------------------------------------------------------------------------------------------
    \begin{frame}{Refrigeradores e Bombas de Calor}\vspace*{-0em}
        \begin{columns}
        \column{0.40\textwidth}
            \begin{itemize}
                \item<1-> O \alert{esquemático} ilustra um \alert{refrigerador genérico};
                    \\[\medskipamount]
                \item<1-> \alert{Sistemas} e \alert{interações energéticas} são identificados;
                    \\[\medskipamount]
                \item<2-> Se \alert{interesse}: \textcolor{green!50!black}{$q_{ent}$}, tem-se um
                    \alert{refrigerador};
					\\[\medskipamount]
                \item<2-> Se \alert{interesse}: \textcolor{green!50!black}{$q_{sai}$}, tem-se uma
                    \alert{bomba de calor}.
            \end{itemize}
        \column{0.60\textwidth}
            \figfigWID{fig/0001-SistemaRefrQsense.pdf}{0.8\textwidth}
		\end{columns}
    \end{frame}
    %-------------------------------------------------------------------------------------------

    %-------------------------------------------------------------------------------------------
    \begin{frame}{Ciclo de Refrigeração por Compressão de Vapor Simples}\vspace*{-0em}
        \figfigHEI{fig/0002-Refr-Vap-ID.pdf}{0.8\textheight}
    \end{frame}
    %-------------------------------------------------------------------------------------------


%-----------------------------------------------------------------------------------------------
\section{Referências}
%-----------------------------------------------------------------------------------------------

    %------------------------------------------------------------------------------------------
    \begin{frame}[allowframebreaks]{Referências -- }
        %% \bibliographystyle{unsrt}
        \bibliographystyle{apalike}
        \setbeamertemplate{bibliography item}{\insertbiblabel}
        \bibliography{bibfile.bib}
    \end{frame}
    %------------------------------------------------------------------------------------------

    % Finishes with stunning image, with credit
    \ArtEndW{pexels-simon-berger-730614}{jpg}{txt}

%-----------------------------------------------------------------------------------------------
\end{document}
%-----------------------------------------------------------------------------------------------

